\section{Introduction}

%%%%%%%%%%%%%%%%%%%%%%%%%%%%%%%%%%%%%%%%%%%%%%%%%%%%%%%%%%%%%%%%%%%%%%%%%%%%%%%%%%%%%%%%%%%%%%%%%%%%%%%%%%%%%%
\subsection{Artistic Nature of Computer Programs}

The study of computers and computations is a science. However, computer programming is an art. To understand
the difference between art and science we can read the following discussion on the subject by John Stuart
Mill:

\begin{quote}
Several sciences are often necessary to form the groundwork of a single art. Such is the complication of human
affairs, that to enable one thing to be done, it is often requisite to know the nature and properties of many
things. ... Art in general consists of the truths of Science, arranged in the most convenient order for
practice, instead of the order which is the most convenient for thought. Science and its truths so as to
enable us to take groups arranges in at one view as much as possible of the general order of the universe. Art
... brings together from parts of the field of science most remote from one another, the truths relating to
the production of the different and heterogeneous conditions necessary to each effect which the exigencies of
practical life require\cite{knuth1992literate}.
\end{quote}

A programmer writes a sequence of instructions to be given the computer so that a desired result can be
achived or `computed by the computer'. This sequence of instructions are what we call a program.
He\footnote{Read `He' as `He/she' throughout this article} has to skillfully use their knowledge about
different subject, then learn some more, and then write the program step by step. There is lots of testing and
debugging involved. There is no well defined way to writing a program. Hence, it is an art. \\

The recent advancements in the field of machine learning have meant that large language models are able to put
together a crappy code to do some stuff, it may very well be possible that we do indeed make computer
programming a science. However, that would not take away the joy of writing programs. AI systems can play
Chess and Go much better than most of the humans on this entire planet, but we still play these games because
they are fun. Same goes for painting and photography --- we can tell a system like DALL-E to cook up a
painting or a drone to take the photograph of forest. However, humans will still continue to paint and click
photographs just because doing so gives us pleasure.

%%%%%%%%%%%%%%%%%%%%%%%%%%%%%%%%%%%%%%%%%%%%%%%%%%%%%%%%%%%%%%%%%%%%%%%%%%%%%%%%%%%%%%%%%%%%%%%%%%%%%%%%%%%%%%
\subsection{Structured Programming Methodology and Literate Programming}

Structured programming methodology arose as a way to introduce rigor into programming. It involves
constructing a program using three basic constructs --- sequential statements (which occur and are executed
one after the other), a conditional statement (\tt if\rm/\tt else\rm), and iterative statments (\tt for \rm
and \tt while \rm loops). Using these three basic constructs we can build simple blocks and functions which
work together to produce the result we desire.\\

Though structured programming allows us to write programs with some discipline, the act of writing a program
is still an art. Some later parts of a program might have been written before the parts which are executed at
the entry point. Moving from $A$ to $B$ in a program might have involved several complex steps, which a bunch
of comments won't explain well. When we see a program, we essentially see a structure which was made out of a
\it web \rm of ideas. This is the reason why we, programmers, have trouble understanding certain pieces of
code which we had written a few months back. When we wrote them, we were in the \it zone \rm and were able to
craft an edifice out of scattered ideas. Once we are done with a code, we are out of that zone. It takes
considerable effort to go from the edifice to the \it web \rm of ideas which created it so that we can
understand the paths to various rooms in the building.\\

We can remedy this by documenting our code. Every good programmer knows the utility of a good documentation.
However, almost every programmer hates to write documentation for their code. Why? Why would we hate to write
something which will only help us and others in understanding something?\\

A major reason for this general dislike is the distractive nature of writing a documentation. Even with all
its utility, it can become a distraction in a world where releases of a code base containing new features and
bug fixes are to be shipped out at a very rapid pace.\\

Also, a documentation might not end up explaining matters clearly sometimes. After all it is explaining the
structure and not the ideas which lead to the creation of the structure. This documentation will be helpful
for someone who wishes to use a library in their project, but it won't be of a great help for those who wish
to understand the underlying construction of the structure. (Of course, you can explain the ideas behind a
piece of code in a traditional documentation, but then why not write literate programs and not break your \it
flow \rm).\\

To remedy the deficiencies of `traditional' programming and documentation, Donald Knuth came up with the
concept of \it Literate Programming\rm  \cite{knuth1984literate}. In literate programming, you write
structured programs. You essentially write an essay which contains `chunks' of the source code present
throughout it in an unstructured yet \it natural \rm way. The prose explains the motivation and working of
chunks. The source code inside a chunk can be expanded by referencing other chunks from inside it or by simply
expanding its definition in the later part of the essay. We get the full structured program by combining the
chunks together, a process which Knuth calls \it tangling \rm as we are producing a tangled up piece out of
web of ideas, something which a computer can understand easily but some person not involved in the development
might have trouble understanding. To produce a documentatiom, you pass the essay through a markup processing
engine like \TeX \ or \LaTeX \ to produce a \tt DVI \rm or \tt PDF \rm file. This is what Knuth calls \it
weaving \rm the web --- you weave the web to produce a beautifully typeset documentation which explains your
code and its API, both for programmers and for users.\\

The original literate programming system was \tt WEB \rm which was used by Knuth when writing \TeX \
typesetting engine. It used a program called \tt tangle \rm which took \tt .web \rm files and output the
source code for \TeX engine written in PASCAL. Another program \tt weave \rm was used to produce the
documentation in \tt DVI \rm format. This documentation was later made available as a book titled \it \TeX :
The Program\rm. The book is available at \tt CTAN \rm, and you can still build both \TeX , the typesetting
engine, and its documentation \cite{knuth1984tex}. \tt WEB \rm could essentially be thought of as a bunch of
macros which were built up on top of \TeX. A version of \tt WEB \rm for C programs called \tt CWEB \rm was
developed by Donald Knuth and Silvio Levy and still used to this day by Knuth and others 
\cite{knuth1992literate}.\\

Sometime later, Norman Ramsey came up with \tt noweb \rm literate programming system 
\cite{ramsey1994literate}. It was an improvement over \tt WEB \rm and \tt CWEB \rm thanks to its simplicity
which arose from the simple markup it used to enclose and refer to code chunks. Code chunks were enclosed
within \texttt{<<Name of the chunk>>=} and \texttt{@}. You could refer to another chunk from inside a chunk by
enclosing the name of the reference inside double angles.

The concept of literate programming is best explained with the help of an example. Say, I want to print the
first $n$ terms of the Fibonacci sequence. My language of choice is C, and here is the code:

\begin{vcode}
//fibonacci.c: Starts at line 23 in web/examples.web
#include <stdio.h>
#include <stdlib.h>
#include <string.h>

int main(int argc, char **argv)
{
    int a, b, n;
    a = -1;
    b = 1;
    if (argc < 2) {
        fprintf(stderr, "Please provide `n' as a command line argument\n");
        exit(1);
    }
    else {
        n = atoi(argv[1]);
    }
    if (n < 0) {
        fprintf(stderr, "n >= 0\n");
        exit(1);
    }
    for (int i = 0; i < n; ++i) {
        fprintf(stdout, "%d ", a + b);
        b = b + a;
        a = b - a;
    }
    fprintf(stdout, "Done\n");
    return 0;
}
\end{vcode}

There is a high chance that you would have skipped the code above. I myself do that!! A literate version of
the program could have been written using noweb markup-style as follows:

\begin{vcode}
<<Fibonacci sequence generator>>=
#include <stdio.h>
#include <stdlib.h>
#include <string.h>

int main(int argc, char **argv)
{
    <<Initialize variables>>
    <<Generate Fibonacci sequence till n terms>>
    fprintf(stdout, "Done\n");
    return 0;
}
\end{vcode}
\verb|@|\\

This gives us the big picture stuff in the fewest of lines. The programmer can then proceed to explain the
working of chunks, for example \coderef{Initialize variables}, and the chunks which make them up. Such a
documentation wouldn't \it repel \rm attention like the plain verbose code did. Even if we were to scatter
some comments through the source code, it (verbatim code) would have still repelled us a bit.\\

\noindent Why?\\

Think about it this way: most humans tend to learn better when they converse with someone, rather than from a
solo fight with the printed letter. The huge popularity of YouTube lectures is an example of that. In a
similar manner, when a programmer writes a program in literate style, he tends to ruminate and explain the
idea behind a piece of code. And these ideas will often come in a random order. This style of presentation, a
short lecture on a piece of code, would be much more engaging than its traditional counterpart.

%%%%%%%%%%%%%%%%%%%%%%%%%%%%%%%%%%%%%%%%%%%%%%%%%%%%%%%%%%%%%%%%%%%%%%%%%%%%%%%%%%%%%%%%%%%%%%%%%%%%%%%%%%%%%%
\subsection{Did Literate Programming take off?}

In a way, yes. We have got Jupyter Notebooks and Mathematica, which provide \it Notebook \rm interfaces where
you write your code like you usually do but intersperse it with chunks of markdown text explaining it. It has
its merits --- scientific collaboration is one of them, which led to the code written in Jupyter notebooks
being used to detect gravitational waves. However, notebooks aren't full literate programming systems as they
do not allow you to place chunks as per your wish or refer to them.\\

The work most cited by fanatics in support of literate programming is \it Physically Based Rendering\rm, a
book on rendering of images in a way that imitates optics in real world, is actually a big literate program
which explains the \tt pbrt \rm library. Its authors won an Academy award as their work was used in movies to
render objects more realistically \cite{pharr2023physically}. Another notable literate program is \it
Axiom\rm, an open-source Computer Algebra System (CAS) maintained by Tim Dally in literate programming style.
He and others have also written a literate version of \it Clojure\rm.\\

Except for a few notable examples, full literate programming systems took off only a bit. Which is sad knowing that had this approach to
programming taken off completely, then programming would have been made more accessible to people and the
quality of code would have improved drastically. Someone on ycombinator said (roughly, I don't remember the
exact words), ``Literate programming requires discipline. It is not suited for world where software
releases are done at such a fast pace.'' --- a statement which I found to be true once I started working in a
corporate. Just think how great it would be if big corporations wrote their code in literate style ---
freelance programmers will be able to understand the codebase and non-technical people (like managers) could
get the big picture of what a program is doing. But... transforming a huge code base into literate style is a
would require too much time and effort, which I don't think any big corporation would be interested in doing.

%%%%%%%%%%%%%%%%%%%%%%%%%%%%%%%%%%%%%%%%%%%%%%%%%%%%%%%%%%%%%%%%%%%%%%%%%%%%%%%%%%%%%%%%%%%%%%%%%%%%%%%%%%%%%%
\subsection{Why then another literate programming tool?}

I could have used \tt noweb \rm but I suffer from NIH - `Not Invented Here' syndrome. The \tt -L \rm in \tt
notangle \rm didn't seem to respect the indentation of reference in parent chunk which is not good for Python
code. Its weaved output was largely unreadable for me. And the hacker guide for it was somewhat confusing ---
I could never figure out how I can change styles or fix the operation of \tt -L \rm
flag \cite{ramsey1992noweb}.\\

So, I decided to write LitCode. It consists of four programs:
\begin{enumerate}
    \item \tt linit \rm: Generates \tt litcode.sty \rm file used for styling and referencing chunks.
    \item \tt ldump \rm: Prints JSON containing chunks present in files passed to it as command line
    arguments.
    \item \tt ltangle \rm: Using the JSON generated by \tt ldump \rm it expands a chunk/reference which has
    been passed to it as command line argument.
    \item \tt lweave \rm: Produces \LaTeX equivalent of files passed to it as command line arguments. It
    transpiles the markup characters used by LitCode to their \LaTeX counterparts.
\end{enumerate}

No one is going to probably use this tool except for me. Given my good experiences with literate programming I
feel that I will be using this system for my life!! So without any more nonsense, lets write the source code
for LitCode.\\

\textbf{Note:} Some code is present in \tt bootstrap/ \rm which is used to build the current version of
programs making up LitCode system. Those pieces are ran first (as you will see in the \tt Makefile \rm in the
parent directory), and then the newly built LitCode is used to tangle and weave the \tt web \rm files.\section{Source Code}

%%%%%%%%%%%%%%%%%%%%%%%%%%%%%%%%%%%%%%%%%%%%%%%%%%%%%%%%%%%%%%%%%%%%%%%%%%%%%%%%%%%%%%%%%%%%%%%%%%%%%%%%%%%%%%
\subsection{Styling file}

I will begin by defining the styling file which will be used by \LaTeX to style the chunks and references.
This makes up \tt litcode \rm package. I will be honest --- I am not that good in \TeX and \LaTeX. Initially I
took lots of help from ChatGPT to cook \tt litcode.sty \rm file. Later on, I decided to keep things simple,
did some additions and deletions, and ended up with the current version of \tt litcode.sty \rm file.

\begin{code}[litcode.sty]
@<$\langle$\textit{litcode.sty}$\rangle\equiv$@>
@<\coderef{Packages required by LitCode}@>
@<\coderef{Variables in LitCode}@>
@<\coderef{hyperref and url style setup}@>
@<\coderef{code environment}@>
@<\coderef{oldcode environment}@>
@<\coderef{vcode environment}@>
@<\coderef{Caption and Chunk Listing format}@>
\end{code}

After lots of hits and trials I ended up with the following list of packages which will be used by LitCode. A
more clever user of \LaTeX \thinspace can maybe add and delete a few packages from this chunk.

\begin{code}[Packages required by LitCode]
@<$\langle$\textit{Packages required by LitCode}$\rangle\equiv$@>
\RequirePackage[T1]{fontenc}
\usepackage{amsmath}
\RequirePackage{textcomp}
%\RequirePackage{inconsolata}
\RequirePackage{calc}
\RequirePackage{caption}
\RequirePackage{hyperref}
\RequirePackage{cleveref}

\RequirePackage{fancyvrb}
\RequirePackage{listings}
\RequirePackage{xcolor}

\RequirePackage{xparse}
\RequirePackage{subcaption}
\end{code}

Let me first tell you about the environments. There are three flavours of environments in LitCode, all of them
being derived from \tt lstlisting \rm environment:

\begin{enumerate}
    \item \tt code \rm: Used when defining a new chunk. Caption and labels will be set for this environment.
    Captions will be invisble so that no weird indentations or spaces are there between normal text and
    listing. \verb|@<| and \verb|@>| are used to escape inside this environment so that we can refer to other
    chunks. 
    \item \tt oldcode \rm: Same as \tt code \rm but used when extending the definition of a previously defined
    chunk. No caption and label is defined for this environment. 
    \item \tt vcode \rm: Used for writing a purely verbatim piece of code. No captions, no labels, no escape characters.
\end{enumerate}

For these environments, I have defined a few variables which will allow us to set listing font, background
color, margin and textwidth, although the latter two aren't used at all (they are just relics of a once
complicated \tt .sty \rm file that I had). 

\begin{code}[Variables in LitCode]
@<$\langle$\textit{Variables in LitCode}$\rangle\equiv$@>
\def\litcode{{LitCode}}
\def\linewidth{\textwidth}

\def\xleftmargin{10pt}
\def\xrightmargin{0pt}

\def\litcodefont{\fontfamily{SourceCodePro-TLF}\small}
\def\codebg{yellow!15}
\def\vcodebg{green!20}

\end{code}

I also wanted to enable `jumps' whenever I would click on reference + I wanted these references to be colored
differently. Thats why I have set \tt colorlinks \rm as \tt true \rm and have set both \tt linkcolor \rm and
\verb|\urlcolor| \rm to \tt red \rm. 

\begin{code}[hyperref and url style setup]
@<$\langle$\textit{hyperref and url style setup}$\rangle\equiv$@>
\hypersetup{
    colorlinks=true,
    linkcolor=red,
    urlcolor=red,
    citecolor=red
}

\end{code}

Thanks to the line below, if you are using a URL, it will be put in the same style as the surrounding text.
\begin{oldcode}
@<$\langle$\textit{hyperref and url style setup}$\rangle \thinspace +\!\!\equiv$@>
\urlstyle{same}
\end{oldcode}

Now come the interesting chunks, which troubled me a lot --- definition of custom environments, caption and
chunk listing style. I will begin with the queen --- the \tt code \rm environment.

\begin{code}[code environment]
@<$\langle$\textit{code environment}$\rangle\equiv$@>
% Define code environment: Used when defining a chunk for the first time
\lstnewenvironment{code}[1][]{
    \lstset{
\end{code}

\coderef{code environment} takes in one argument which is used to set both the caption and label of the environment.
Note how I am using \verb|\protect\detokenize{}| to make sure the formatting is not messed up. Do note that
captions will be invisible so that we can have captions and code chunk indented by the same level. Also `fake
captions' will be inserted to by \tt lweave \rm with proper formatting so that chunk definitions and
extensions can be distinguished from each other.

\begin{oldcode}
@<$\langle$\textit{code environment}$\rangle \thinspace +\!\!\equiv$@>
        caption={\protect\detokenize{#1}},
        label={#1},
\end{oldcode}

We further extend \coderef{code environment} by setting \tt columns \rm to \tt fullflexible \rm so that the letters
are nicely typeset (unlike in \tt fixed \rm column width). I have set the font to \verb|\ttfamily| --- you can
change it to something else. \tt keepspaces \rm, and \tt breaklines \rm are used to make sure the indentation
is not changed and the lines are wrapped if they go over margins, respectively. \tt upquote = true \rm ensure
that double quotes are formatted as simple upquotes only (and not as `` and ''). \tt background \rm is used to
set the background color via \verb|\codebg| variable: set it to any color you wish in the preamble.

\begin{oldcode}
@<$\langle$\textit{code environment}$\rangle \thinspace +\!\!\equiv$@>
        columns=fullflexible,
        basicstyle=\litcodefont,
        keepspaces=true,
        breaklines=true,
        upquote=true,
        backgroundcolor=\color{\codebg},
\end{oldcode}

The interesting part is the \tt escapeinside \rm variable, which is used to set the characters which will be
used to escape inside the \tt lstlisting \rm environment. I have set these variables to \verb|@<| and
\verb|@>|. To print these variables inside the environment, they need to be escaped themselves. Example:

\begin{vcode}
@< --> gets transpiled to --> @<@<@>
@> --> gets transpiled to --> @<@>@>
\end{vcode}

\begin{oldcode}
@<$\langle$\textit{code environment}$\rangle \thinspace +\!\!\equiv$@>
        escapeinside={@<@<@>}{@<@>@>}
    }
}{}

\end{oldcode}

We can define \tt oldcode \rm and \tt vcode \rm environments in a similar fashion with a few minor changes: in
\tt oldcode \rm we will get rid of caption and label, and in \tt vcode \rm we will get rid of escape
characters.

\begin{code}[oldcode environment]
@<$\langle$\textit{oldcode environment}$\rangle\equiv$@>
% Define oldcode environment: To be used when extending an already defined chunk
\lstnewenvironment{oldcode}[1][]{
    \lstset{
        columns=fullflexible,
        basicstyle=\litcodefont,
        keepspaces=true,
        breaklines=true,
        upquote=true,
        backgroundcolor=\color{\codebg},
        escapeinside={@<@<@>}{@<@>@>}
    }
}{}
\end{code}

\begin{code}[vcode environment]
@<$\langle$\textit{vcode environment}$\rangle\equiv$@>
% Verbatim code environment: Takes no captions, has no labels and won't get included in the listing. 
% Purely verbatim!!
\lstnewenvironment{vcode}[1][]{
    \lstset{
        columns=fullflexible,
        basicstyle=\litcodefont,
        keepspaces=true,
        breaklines=true,
        upquote=true,
        backgroundcolor=\color{\vcodebg},
    }
}{}
\end{code}

We end this \tt .sty \rm file by defining the format in which we want to display captions in references and in
listings. I have followed the style which had been used by Knuth and Ramsey:

\begin{itemize}
    \item $\langle$ name of the chunk $\rangle \equiv$: Used when defining a chunk for the first time
    \item $\langle$ name of the chunk $\rangle \thinspace +\!\!\equiv$: Used when extending the definition of
    a chunk
    \item $\langle$ name of the chunk $\rangle$: Used when referencing a chunk
\end{itemize}

The first and third type of formatting is defined in \tt litcode.sty \rm itself. One can cook up a \LaTeX code
which checks whether a chunk has been previously defined or not and accordingly format its captioning. I am
not that clever. This decision (checking the \it freshness \rm of the chunk and formatting its caption) is
done by \tt lweave \rm. 

\begin{code}[Caption and Chunk Listing format]
@<$\langle$\textit{Caption and Chunk Listing format}$\rangle\equiv$@>
\DeclareCaptionFormat{codecaptionformat}
{
    \phantom{$\langle\textit{#3}\rangle\!\!\equiv$}
}

\captionsetup[lstlisting]{
    format=codecaptionformat,
    justification=raggedright,
    singlelinecheck=off,
    skip=-\baselineskip
}

\makeatletter
\renewcommand\l@lstlisting[2]{\@dottedtocline{1}{0em}{2.3em}{\textit{#1}}{#2}}
\makeatother

\newcommand{\coderef}[1]{$\langle$\textit{\nameref{#1}}$\rangle$}
%\newcommand{\coderef}[1]{$\langle$\textit{\nameref{#1}} \rm(p.~\pageref{#1})$\rangle$}
\end{code}

Note how in \verb|\DeclareCaptionFormat| I have set the caption under \verb|\phantom| command so that they are
invisble and do not mess up the indentation. In \verb|\captionsetup| I have setup \tt skip \rm to
\verb|-\baselineskip| so that the vertical space between a chunk and lines of normal text preceding it is not
huge.

%%%%%%%%%%%%%%%%%%%%%%%%%%%%%%%%%%%%%%%%%%%%%%%%%%%%%%%%%%%%%%%%%%%%%%%%%%%%%%%%%%%%%%%%%%%%%%%%%%%%%%%%%%%%%%
\subsection{linit.py}

\tt linit.py \rm or simply \tt linit \rm, will be responsible for writing \tt litcode.sty \rm to a directory
which is a passed to it as a command line argument. Of course of course, I can upload this styling file CTAN,
but I do not think that this is anything ground-breaking, so I will keep it with myself.\\

The first line of \tt linit.py \rm and all other Python scripts will begin with a shebang (\verb|#!|)
followed by the path to the binary of python3 interpreter. All these scripts will be set as executables (\tt
chmod +x \rm) and then moved to \tt ~/local/bin/ \rm so that we can run them from anywhere in the system
without explicitely calling Python to execute them everytime.

\begin{code}[linit.py]
@<$\langle$\textit{linit.py}$\rangle\equiv$@>
#!/usr/bin/python3
import os
import sys
import time
import subprocess as sbpr

\end{code}

The entire content of \tt litcode.sty \rm will be present inside \tt linit.py \rm as a multi-line raw string.
Why a raw string? This is because this \tt .sty \rm file makes extensive use of backslash \verb|\|. I do not
want to make things messy by escaping every backslash... Better to just use a raw string. :-)

\begin{oldcode}
@<$\langle$\textit{linit.py}$\rangle \thinspace +\!\!\equiv$@>
litcode_sty_content = r'''
@<\coderef{litcode.sty}@>
'''

\end{oldcode}

\tt linit.py \rm will accept only a single command line argument which would be the name of the directory
where \tt litcode.sty \rm is to be saved. The `entry-point' for our script will be the \tt main \rm function,
where the command line arguments will be processed to get this directory name. If none is supplied, then the
program should exit with an error code.

\begin{oldcode}
@<$\langle$\textit{linit.py}$\rangle \thinspace +\!\!\equiv$@>
def main():
    global litcode_sty_content
    if len(sys.argv) <= 1:
        print('Please supply a directory name')
        exit(1)
    directory = sys.argv[1]
\end{oldcode}

We can also ensure that the directory is made using \tt mkdir -p \rm. It won't do any harm to the directory's
contents if it already exists.

\begin{oldcode}
@<$\langle$\textit{linit.py}$\rangle \thinspace +\!\!\equiv$@>
    command = f'mkdir -p {directory}'
    os.system(command)
\end{oldcode}

Finally, we write the raw string to \tt litcode.sty \rm placed under the directory supplied as a command line
argument. We need the full path to \tt litcode.sty \rm which can easily be made by appending the name of the
directory (relative to the current directory/position in the terminal), system path separator (which we get by
\tt os.sep()\rm) and the string \verb|'litcode.sty'|. I have just used a formatted string, although these
strings could also have been appended by \tt + \rm operator.

\begin{oldcode}
@<$\langle$\textit{linit.py}$\rangle \thinspace +\!\!\equiv$@>
    full_path_to_sty = f'{directory}{os.sep}litcode.sty'
    with open(full_path_to_sty, 'w') as fp:
        fp.writelines(litcode_sty_content)

main()
\end{oldcode}

%%%%%%%%%%%%%%%%%%%%%%%%%%%%%%%%%%%%%%%%%%%%%%%%%%%%%%%%%%%%%%%%%%%%%%%%%%%%%%%%%%%%%%%%%%%%%%%%%%%%%%%%%%%%%%
\subsection{ldump.py}

Tangling a chunk involves recursive expansion of all the references that are there in the chunk. The
indentation of a reference must be appended while printing all the lines of the chunk it refers to. We also
need to add some markers which tell a programmer from where a piece of code lies in the WEB files. We also
need to ensure that proper whitespace conversion is done --- this is especially useful in the case of \tt
Makefile \rm which accepts only tabs as valid indentation characters.\\

Most of the traditional literate programming tools tangle a chunk by going through a file, making a dictionary
of chunks and then expanding the required chunk recursively. What would happen if we need to tangle 2 chunks
from the same file. Some computational power will be expended in setting up the dictionary twice. To prevent
this `waste of computational power' (although in the 21st century, hardly anyone cares about it for such
simple text processing job), I have written \tt ldump.py \rm. It will make a dicitionary of chunks which are
present in the files which have been supplied to it as command line arguments. Then this dictionary will be
printed and can be saved in a JSON file using the clobber operator, \tt > \rm. This JSON will be used by \tt
ltangle \rm when tangling a piece of code. \\

\textbf{NOTE:} If a chunk is defined in multiple files and have no relation to each other, give them different
names. This will prevent \tt ltangle \rm and \tt lweave \rm from considering them as related to each other.\\

We begin with the shebang line specifying the path of the python3 and import \tt os\rm, \tt sys\rm, \tt
json \rm and \tt re \rm modules.

\begin{code}[ldump.py]
@<$\langle$\textit{ldump.py}$\rangle\equiv$@>
#!/usr/bin/python3
import os
import sys
import re
import json

\end{code}

A pseudocode for \tt ldump.py \rm is presented in Algorithm~\ref{algo:ldump-working}:

\begin{algorithm}

Check if command line args have been supplied. \linebreak
If yes: then start reading each file line by line \linebreak
If not: then throw an error\\

Go through the string and look lines satisfying the regex: \linebreak
    \verb|(?<=^<<).+?(?=>>=\s*$)|\\

Check if the chunk exists or not. If it is not their in the dictionary, then enter it. Else, we will append
the lines in this chunk to the already existing lines in the dictionary.\\

Add some metadata about the position of chunk boundaries in the web files.\\

JSONify the dictionary and print it. Writing to the file will be taken care of by \tt > \rm operator.

\caption{Pseudocode for \texttt{ldump.py}}
\label{algo:ldump-working}
\end{algorithm}

We will be having a function \tt dumper \rm which will take in a list of files, read them line by line and
create a dictionary of chunks, \tt chunkDict \rm. This dictionary will then be converted into a pretty-printed
JSON string using \tt dumps \rm method in \tt json \rm package.

\begin{oldcode}
@<$\langle$\textit{ldump.py}$\rangle \thinspace +\!\!\equiv$@>
def dumper(file_list):
    chunkDict = dict()
    for f in file_list:
        @<\coderef{Fill chunkDict}@>
    chunkJSON = json.dumps(chunkDict, indent = 4)
    return chunkJSON

\end{oldcode}

Our entry point into the program will be \tt main() \rm which will see if the files have been supplied as
command line args or not. If they exist, we will send this list to the \tt dumper \rm. The JSON string that it
gets from the \tt dumper \rm will be printed on the screen.

\begin{oldcode}
@<$\langle$\textit{ldump.py}$\rangle \thinspace +\!\!\equiv$@>
def main():
    if len(sys.argv) <= 1:
        print('No files provided')
        exit(1)
    file_list = sys.argv[1:]
    json_str = dumper(file_list)
    print(json_str)

main()
\end{oldcode}

We now have to write some code which will write \tt chunkDict \rm. When reading a file, we need to keep track
of the following:

\begin{enumerate}
    \item Line number of the current line being read.
    \item Are we in a chunk or not?
        \begin{itemize}
            \item If not, then does this line defines a chunk? Is it really new or an old one?
            \item If yes, then does this line refers to another chunk? Is there a \verb|@| in this line?
            (which marks the end of a chunk's definition)
        \end{itemize}
\end{enumerate}

We keep track of these conditions with the help of \verb|line_count|, \tt inChunk\rm, \tt
encounteredReference\rm, and \tt encounteredStopChar\rm variables.

\begin{code}[Fill chunkDict]
@<$\langle$\textit{Fill chunkDict}$\rangle\equiv$@>
line_count = 0
inChunk = False
encounteredReference = False
encounteredStopChar = False

\end{code}

We will read all the lines of a file and store them in a list (\tt lines \rm) in one shot using \tt readlines
\rm method. \tt chunkLines \rm will store the lines of a chunk as and when we encounter. It will be off-loaded
into \tt chunkDict \rm.

\begin{oldcode}
@<$\langle$\textit{Fill chunkDict}$\rangle \thinspace +\!\!\equiv$@>
with open(f, 'r') as fp:
    lines = fp.readlines()
chunkLines = []

\end{oldcode}

Using a \tt for \rm loop I am now going to iterate over \tt lines \rm. The first step in each iteration would
be to increment \verb|line_count|.

\begin{oldcode}
@<$\langle$\textit{Fill chunkDict}$\rangle \thinspace +\!\!\equiv$@>
for line in lines:
    line_count += 1
    @<\coderef{Fill chunkDict: Not in a chunk}@>
    @<\coderef{Fill chunkDict: In a chunk}@>
\end{oldcode}

There will be two conditions. Either I will be in a chunk or I won't. If I am not in a chunk, \tt inChunk \rm
is set to \tt False \rm, and I have to check whether the current line is the beginning of a chunk definition
or not.

\begin{code}[Fill chunkDict: Not in a chunk]
@<$\langle$\textit{Fill chunkDict: Not in a chunk}$\rangle\equiv$@>
if not inChunk:
    chunkname = re.findall(r'(?<=^<<).+?(?=>>=\s*$)', line.strip())
    if chunkname:
        inChunk = True

\end{code}

I now need to check whether the chunk has been previously defined or not. This can be easily done with the use
of \tt in \rm operator. Remember that \tt re.findall() \rm returns a list, so I need to set \tt chunkname \rm
to the first element of the list that has been returned by \tt findall()\rm.

\begin{oldcode}
@<$\langle$\textit{Fill chunkDict: Not in a chunk}$\rangle \thinspace +\!\!\equiv$@>
        chunkname = chunkname[0]
        if chunkname not in chunkDict:
            chunkDict[chunkname] = []
            chunkLines = []
            chunkLines.append(f'@->{chunkname}: Starts at line {line_count} in {f}\n')
        else:
            chunkLines = chunkDict[chunkname]
            chunkLines.append(f'@->{chunkname}: Extended at line {line_count} in {f}\n')

\end{oldcode}

Note how I am introducing metadata beginning with \verb|@->| at the beginning of the chunk. \\

I have to now deal with the condition in which I am inside a chunk. When inside a chunk, I need to check
whether the current line contains a reference or a stop-character (signaling end of the chunk), and
accordingly insert metadata at chunk boundaries.

\begin{code}[Fill chunkDict: In a chunk]
@<$\langle$\textit{Fill chunkDict: In a chunk}$\rangle\equiv$@>
else:
    encounteredStopChar = (line.strip() == '@')
    if encounteredStopChar:
        chunkLines.append(f'@->{chunkname}: Ends at {line_count} in {f}\n')
        chunkDict[chunkname] = chunkLines
        inChunk = False
        continue
    else:
        encounteredReference = re.findall(r'(?<=^<<).+?(?=>>\s*$)', line.strip())
        if encounteredReference:
            chunkLines.append(f'@->{chunkname}: Reference at line {line_count} in {f}\n')
            chunkLines.append(line)
            chunkLines.append(f'@->{chunkname}: Continues from line {line_count + 1} in {f}\n')
        else:
            chunkLines.append(line)
\end{code}

%%%%%%%%%%%%%%%%%%%%%%%%%%%%%%%%%%%%%%%%%%%%%%%%%%%%%%%%%%%%%%%%%%%%%%%%%%%%%%%%%%%%%%%%%%%%%%%%%%%%%%%%%%%%%%
\subsection{ltangle.py}

Much of the heavy-lifting of collecting all chunks and putting them in a dictionary + putting some metadata
for debugging purposes had already been done \tt ldump\rm. We now have to write \tt ltangle.py \rm which
will just need to expand a chunk/reference which has been passed to it as a command line argument. Let's make
a list of thing which will be passed to \tt ltangle \rm as command line arguments, along with relevant
switches. 

\begin{enumerate}
    \item \tt -i dump/file.json\rm: \tt ltangle \rm will be using the dump file generated by \tt ldump \rm to
    read chunks and expand them recursively.
    \item \tt -R name/of/the/chunk\rm: The name of the chunk we want to expand. 
    \item \tt -cc comment/line/starting/character\rm: Character which is to be appended before metadata lines
    so as to comment them out in the tangled source code. For example, \verb|//| would be used for commenting
    lines in a C or C++ code. If no comment-char is provided, no metadata line will be included.
    \item \tt -ut number/of/spaces/to/convert\rm: Convert specified number of spaces to tabs. Particularly
    useful for systems which accept only tabs (or spaces) for indentation, example: Makefile.
    \item \tt -us number/of/tabs/to/convert\rm: Convert tabs to specified number of spaces.
\end{enumerate}

Like before, we begin our code with a shebang. We import the same modules as before. 

\begin{code}[ltangle.py]
@<$\langle$\textit{ltangle.py}$\rangle\equiv$@>
#!/usr/bin/python3
import os
import sys
import re
import json

\end{code}

I am setting \tt chunkDict \rm as a global variable. This eliminates the need to pass it as an argument during
recursive expansion of the reference.

\begin{oldcode}
@<$\langle$\textit{ltangle.py}$\rangle \thinspace +\!\!\equiv$@>
chunkDict = dict()

\end{oldcode}

\tt expandref \rm is the queen of this program. It takes in a single parameter --- the name of the reference
which is to be expanded recursively.

\begin{oldcode}
@<$\langle$\textit{ltangle.py}$\rangle \thinspace +\!\!\equiv$@>
def expandref(reference):
    @<\coderef{ltangle: expandref}@>

\end{oldcode}

Like before, our entry point will be the \tt main \rm function where we will perform much of the processing of
the command-line arguments.

\begin{oldcode}
@<$\langle$\textit{ltangle.py}$\rangle \thinspace +\!\!\equiv$@>
def main():
    @<\coderef{ltangle: main}@>

main()
\end{oldcode}

So, this would be the core structure of \tt ltangle\rm. Let's start writing \coderef{ltangle: expandref} as that is
the most important part of this program. The expanded reference will be kept in \tt codeLines \rm which will
be returned to \tt main \rm for further processing.\\

\begin{code}[ltangle: expandref]
@<$\langle$\textit{ltangle: expandref}$\rangle\equiv$@>
global chunkDict
codeLines = []
@<\coderef{Expand the reference}@>
return codeLines
\end{code}

When expanding reference, we need to take care of the indentation of the reference. This is especially
important when we are dealing with a language where indentation defines a block of code, like in Python. The
pseudo-code for \coderef{Expand the reference} will be something like this:

\begin{algorithm}
Read a line\\
Check if it contains a reference. \linebreak
\textbf{if} it contains one, call yourself (\texttt{expandref()}) again and expand
upon this reference. Once you get the lines making up this reference, to each line append the whitespace with
which this reference was indented. Then append it to \tt codeLines.\rm\linebreak
\textbf{else} simply append the line to \tt codeLines\rm.

\caption{Pseudocode for \texttt{ltangle.py}}
\label{algo:ltangle-working}
\end{algorithm}

I will be using two variables, \tt encounteredReference \rm and \texttt{indentation} to store the current
state of the function before possible recurrsive call.

\begin{code}[Expand the reference]
@<$\langle$\textit{Expand the reference}$\rangle\equiv$@>
encounteredReference = False
indentation = ''
\end{code}

I pass the name of the reference to \texttt{chunkDict} and in return I get a list of lines making up that
reference. I iterate over them. At each iteration I check whether the line contains a reference or not. If it
doesn't, I simply append it to \texttt{chunkDict} and move ahead with the next iteration. If it does, then as
per \ref{algo:ltangle-working}, I am going to call \tt expandref\rm with this reference and get its expansion.
This will be stored in \tt buffer\rm.

\begin{oldcode}
@<$\langle$\textit{Expand the reference}$\rangle \thinspace +\!\!\equiv$@>
for line in chunkDict[reference]:
    encounteredReference = re.findall(r'(?<=^<<).+?(?=>>\s*$)', line.strip())
    if encounteredReference:
        reference = encounteredReference[0]
        encounteredReference = True
        buffer = expandref(reference)
\end{oldcode}

Now, I extract the whitespace using which we had expanded the reference and append it to each and every line
in the \tt buffer\rm.
\begin{oldcode}
@<$\langle$\textit{Expand the reference}$\rangle \thinspace +\!\!\equiv$@>
        indentation = line[:len(line) - len(line.lstrip())]
        for l in buffer:
            line = indentation + l
            codeLines.append(line)
\end{oldcode}

Under the event no reference has been encountered, I simply append the line to \tt codeLines\rm and carry on
with the next iteration.
\begin{oldcode}
@<$\langle$\textit{Expand the reference}$\rangle \thinspace +\!\!\equiv$@>
    else:
        codeLines.append(line)
\end{oldcode}

Coming to \tt main\rm, I need to do some processing of command-line arguments to understand how the tangled
code has to be formatted.

\begin{code}[ltangle: main]
@<$\langle$\textit{ltangle: main}$\rangle\equiv$@>
global chunkDict
dumpfile = str()
reference = str()
commentchar = None
usetabs = False
whitespace_conv_num = 4
if '-i' in sys.argv:
    i = sys.argv.index('-i')
    dumpfile = sys.argv[i + 1]
else:
    print('Cannot find -i flag')
    exit(1)
if '-R' in sys.argv:
    i = sys.argv.index('-R')
    reference = sys.argv[i + 1]
else:
    print('Cannot find -R flag')
    exit(1)
if '-cc' in sys.argv:
    i = sys.argv.index('-cc')
    commentchar = sys.argv[i + 1]
if '-ut' in sys.argv:
    usetabs = True
    i = sys.argv.index('-ut')
    whitespace_conv_num = int(sys.argv[i + 1])
if '-us' in sys.argv:
    usetabs = False
    i = sys.argv.index('-us')
    whitespace_conv_num = int(sys.argv[i + 1])

\end{code}

If all is good with the arguments that have been passed to us, we are then going to load the dump file, set
\tt chunkDict\rm and then call \tt expandref \rm passing reference as an argument to it.

\begin{oldcode}
@<$\langle$\textit{ltangle: main}$\rangle \thinspace +\!\!\equiv$@>
# This is the portion where we can load the JSON file
chunkDict = dict()
with open(dumpfile, 'r') as fp:
    chunkDict = json.load(fp)

# Now we call the expandref function
codeLines = expandref(reference)
#print(''.join(codeLines))

\end{oldcode}

Now, we have to format the lines --- perform some white-space conversion and comment out metadata lines.

\begin{oldcode}
@<$\langle$\textit{ltangle: main}$\rangle \thinspace +\!\!\equiv$@>
# Now replace '@->' with comment character
lines = []
for line in codeLines:
    @<\coderef{Perform whitespace conversion}@>
    @<\coderef{Comment out metadata lines}@>
    lines.append(line)
codeLines = lines

\end{oldcode}

Let us say the line was originally indented by $s$ spaces, which we wish to replace by $t$ tabs and $u$
spaces, with each tab being equivalent to $n_s$ spaces, then we have to implement the following equations: 

$$t = \left\lfloor \frac{s}{n_s} \right\rfloor$$
$$u = s\ \mathbf{mod}\ n_s$$

\begin{code}[Perform whitespace conversion]
@<$\langle$\textit{Perform whitespace conversion}$\rangle\equiv$@>
# Spaces to Tabs
if usetabs is True:
    indentation_depth = len(line) - len(line.lstrip(' '))
    # This condition is required, otherwise for indentation_depth = 0, all
    # indentation will be eaten away.
    if indentation_depth:
        new_indentation = '\t' * (indentation_depth // whitespace_conv_num)
        new_indentation += ' ' * (indentation_depth % whitespace_conv_num)
        line = new_indentation + line.lstrip()
# Tabs to Spaces
elif usetabs is False:
    indentation_depth = len(line) - len(line.lstrip('\t'))
    if indentation_depth:
        new_indentation = ' ' * (indentation_depth * whitespace_conv_num)
        line = new_indentation + line.lstrip()
\end{code}

Commenting metadata lines is pretty simple. Just check if the line starts with \verb|@->| character. If it
does, replace it with \tt commentchar\rm  if it has been supplied in the command line, else just skip it.

\begin{code}[Comment out metadata lines]
@<$\langle$\textit{Comment out metadata lines}$\rangle\equiv$@>
if line.lstrip().startswith('@->'):
    if commentchar is not None:
        line = line.replace('@->', commentchar, 1)
    else:
        continue

\end{code}

Finally, we are going to print \tt codeLines\rm. The output can be redirected to write to a file using \tt >
\rm operator.

\begin{oldcode}
@<$\langle$\textit{ltangle: main}$\rangle \thinspace +\!\!\equiv$@>
# Now print it
print(''.join(codeLines))

\end{oldcode}

%%%%%%%%%%%%%%%%%%%%%%%%%%%%%%%%%%%%%%%%%%%%%%%%%%%%%%%%%%%%%%%%%%%%%%%%%%%%%%%%%%%%%%%%%%%%%%%%%%%%%%%%%%%%%%
\subsection{lweave.py}

Now we have to write the final program of the LitCode system, \tt lweave.py\rm. Like before, we begin with a
shebang and the necessary imports.

\begin{code}[lweave.py]
@<$\langle$\textit{lweave.py}$\rangle\equiv$@>
#!/usr/bin/python3
import os
import sys
import re

\end{code}

Again, \tt main \rm will be our entry point in this program. \tt lweave \rm should read all the files that
have been passed to it as command-line arguments, club together the lines making them, and transpile them to
valid \LaTeX code. This transpilation will be done by a function named \tt transpiler \rm.

\begin{oldcode}
@<$\langle$\textit{lweave.py}$\rangle \thinspace +\!\!\equiv$@>

def transpiler(lines):
    @<\coderef{transpile from web to latex}@>

def main():
    @<\coderef{lweave: main}@>

main()
\end{oldcode}

First let us get \coderef{lweave: main} out of the way. If no files are provided to it, \tt lweave \rm should fail.

\begin{code}[lweave: main]
@<$\langle$\textit{lweave: main}$\rangle\equiv$@>
if len(sys.argv) < 2:
    print('File names have to be provided as command line arguments')
    exit(1)

\end{code}

If files have been provided to it, lines of each file will be read and stored in a common list, which will
then be sent to \tt transpiler \rm. The transpiler will return another list, containing the transpiled lines.

\begin{oldcode}
@<$\langle$\textit{lweave: main}$\rangle \thinspace +\!\!\equiv$@>
files = sys.argv[1:]
lines = []
transpiled_lines = []
for file in files:
    with open(file, 'r') as fp:
        lines += fp.readlines()
transpiled_lines += transpiler(lines)

\end{oldcode}

These lines will then be joined to form a single string and will be printed on the screen.

\begin{oldcode}
@<$\langle$\textit{lweave: main}$\rangle \thinspace +\!\!\equiv$@>
print(''.join(transpiled_lines))

\end{oldcode}

The transpiler needs to check whether it is reading a normal code line or a line in some listing, i.e., from
\tt code\rm, \tt oldcode\rm, or \tt vcode \rm environment and take actions accordingly. The transpilations it
need to perform in each of the environments are as follows:

\begin{enumerate}
    \item \textbf{Normal mode}
    \begin{itemize}
        \item[] Text within \verb|<<| and \verb|>>| to be extracted. Check if such a reference exists. If it
        exists, format to \verb|\coderef|, else continue without alteration
        \item[] \verb|@<| and \verb|@>| should remain unaltered 
    \end{itemize}
    \item \textbf{Listing mode}
    \begin{itemize}
        \item[] Text within \verb|<<| and \verb|>>=| to be extracted. If \verb|len(text)| is zero, then
        replace by \verb|\begin{vcode}|. If not, then check if the reference has already been defined or not.
        If it is a new one, perform substitution with \verb|\begin{code}| and \verb|<<extracted text>>=|. Else
        with \verb|\begin{oldcode}| and \verb|<<extracted text>> +=|
        \item[] Text within \verb|<<| and \verb|>>| to be extracted. Substitution to be done with
        \verb|@<\coderef{}@>|.
        \item[] \verb|@<| and \verb|@>| to be escaped properly.
    \end{itemize}
\end{enumerate}

Given the transpilations, here are the variables which we will use to keep track of things for us.

\begin{code}[transpile from web to latex]
@<$\langle$\textit{transpile from web to latex}$\rangle\equiv$@>
mode = 'Normal'
chunkList = []
processed_chunkList = []
transpiled_lines = []

\end{code}

We are going to go through all the strings and make a list of chunks present in the files.

\begin{oldcode}
@<$\langle$\textit{transpile from web to latex}$\rangle \thinspace +\!\!\equiv$@>
for line in lines:
    chunkname = re.findall(r'(?<=^<<).+?(?=>>=\s*$)', line.strip())
    if chunkname:
        chunkname = chunkname[0]
        if len(chunkname) != 0:
            chunkList.append(chunkname)

\end{oldcode}

We are again going to iterate over the lines, but this time actually transpile them. 

\begin{oldcode}
@<$\langle$\textit{transpile from web to latex}$\rangle \thinspace +\!\!\equiv$@>
for line in lines:
    if mode == 'Normal':
        @<\coderef{transpiler: Normal mode}@>
    
    elif mode == 'code' or mode == 'oldcode':
        @<\coderef{transpiler: code and oldcode mode}@>
    
    elif mode == 'vcode':
        @<\coderef{transpiler: vcode mode}@>
    
    transpiled_lines.append(line)
return transpiled_lines

\end{oldcode}

Now let us write `rules' of transpiling web to \LaTeX in normal mode. We will first check whether we have
encountered a listing environment or not and if we have, we need to set the mode and add some lines accordingly,

\begin{code}[transpiler: Normal mode]
@<$\langle$\textit{transpiler: Normal mode}$\rangle\equiv$@>
chunkname = re.findall(r'(?<=^<<).+?(?=>>=\s*$)', line.strip())
if chunkname:
    chunkname = chunkname[0]
    if chunkname not in processed_chunkList:
        mode = 'code'
        processed_chunkList.append(chunkname)
        envbegin = '\\begin{code}[' + chunkname + ']\n'
        caption = '@<@<@>$\\langle$\\textit{' + chunkname + '}$\\rangle\\equiv$@<@>@>\n'
    else:
        mode = 'oldcode'
        envbegin = '\\begin{oldcode}\n'
        caption = '@<@<@>$\\langle$\\textit{' + chunkname + '}$\\rangle ' + \
                    '\\thinspace +\\!\\!\\equiv$@<@>@>\n'
    transpiled_lines.append(envbegin)
    transpiled_lines.append(caption)
    continue

\end{code}

\tt code \rm and \tt oldcode \rm will be set if there is some text within the angled brackets. But what if
there is no text? In such case, \verb|<<>>=| will mark the beginning of \tt vcode \rm environment.

\begin{oldcode}
@<$\langle$\textit{transpiler: Normal mode}$\rangle \thinspace +\!\!\equiv$@>
vcodeBegins = (line.strip() == '<<>>=')
if vcodeBegins:
    mode = 'vcode'
    envbegin = '\\begin{vcode}\n'
    transpiled_lines.append(envbegin)
    continue

\end{oldcode}

If none of these environments have started, then we are simply in \it Normal \rm mode. Then \verb|@<| and
\verb|@>| should remain unaltered. However, if we are encountering a reference, and if it defined in the file
in the line, it should be transpiled to \verb|\coderef| (as has been mentioned above).

\begin{oldcode}
@<$\langle$\textit{transpiler: Normal mode}$\rangle \thinspace +\!\!\equiv$@>
reference = re.findall(r'(?<=<<).+?(?=>>.*)', line.strip())
if reference:
    reference = reference[0]
    if reference in chunkList:
        line = line.replace('<<', '\\coderef{')
        line = line.replace('>>', '}')

\end{oldcode}

If we are in either \textbf{code} or \textbf{oldcode} mode, we need to escape \verb|@<| and
\verb|@>| so that they can be printed properly. Also, any references should be transpiled to \verb|\coderef|
and be escaped by the escape characters.

\begin{code}[transpiler: code and oldcode mode]
@<$\langle$\textit{transpiler: code and oldcode mode}$\rangle\equiv$@>
buffer = ''
i = 0
while i < len(line):
    if line[i] == '@':
        if i < len(line) - 1:
            if line[i + 1] == '<' or line[i + 1] == '>':
                buffer += '@<@<@>' + line[i] + line[i + 1] + '@<@>@>'
                i += 2
                continue
    buffer += line[i]
    i += 1
line = buffer

reference = re.findall(r'(?<=^<<).+?(?=>>\s*$)', line.strip())
if reference:
    reference = reference[0]
    if reference in chunkList:
        line = line.replace('<<', '@<@<@>\\coderef{')
        line = line.replace('>>', '}@<@>@>')

\end{code}

We also need to make sure that if we are encountering \verb|@| we need to switch modes immediately.

\begin{oldcode}
@<$\langle$\textit{transpiler: code and oldcode mode}$\rangle \thinspace +\!\!\equiv$@>
if line.strip() == '@':
    line = '\\end{' + mode + '}\n'
    mode = 'Normal'

\end{oldcode}

The code for transpiling \texttt{vcode} is relatively simple. Everything has to be printed out verbatim. We
just need to check whether we have encountered stop-char or not.

\begin{code}[transpiler: vcode mode]
@<$\langle$\textit{transpiler: vcode mode}$\rangle\equiv$@>
if line.strip() == '@':
    mode = 'Normal'
    line = '\\end{vcode}\n'

\end{code}
\section{Examples}
\subsection{Fibonacci Sequence}
The C code presented in this section is responsible for printing the Fibonacci sequence.\\

I begin by writing the \tt Makefile \rm for building the program.
\begin{code}[Makefile]
@<$\langle$\textit{Makefile}$\rangle\equiv$@>
CC := gcc -std=C99 -O0

runfib: fibonacci.exe
    ./fibonacci.exe 30

fibonacci.exe: fibonacci.o
    $(CC) fibonacci.o -o fibonacci.exe

fibonacci.o: fibonacci.c
    $(CC) -c fibonacci.c -o fibonacci.o
\end{code}

Now comes the fun part. Our code will have two parts: \coderef{Initialize variables} will be responsible for
initializing the variables which will then be used by \coderef{Generate Fibonacci sequence till n terms} to generate
the sequence till some $n$ number of terms. $n$ will be provided by the user as a command line argument.

\begin{code}[fibonacci.c]
@<$\langle$\textit{fibonacci.c}$\rangle\equiv$@>
#include <stdio.h>
#include <stdlib.h>
#include <string.h>

int main(int argc, char **argv)
{
    @<\coderef{Initialize variables}@>
    @<\coderef{Generate Fibonacci sequence till n terms}@>
    fprintf(stdout, "Done\n");
    return 0;
}
\end{code}

To generate Fibonacci sequence, we need three numbers: $a$, $b$, and $n$. $a$ and $b$ will be used to generate
the terms of the sequence till $n$ terms. $n$ will be a command line argument. We will need to check whether
$n$ has been passed or not.

\begin{code}[Initialize variables]
@<$\langle$\textit{Initialize variables}$\rangle\equiv$@>
int a, b, n;
a = -1;
b = 1;
@<\coderef{Check cmd args and Initialize n}@>
\end{code}

\begin{code}[Check cmd args and Initialize n]
@<$\langle$\textit{Check cmd args and Initialize n}$\rangle\equiv$@>
if (argc < 2) {
    fprintf(stderr, "Please provide `n' as a command line argument\n");
    exit(1);
}
else {
    n = atoi(argv[1]);
}
\end{code}

We also need to make sure that $n$ is not less than zero. So, we will extend the definition of <<Check cmd
args and Initialize n>>.

\begin{oldcode}
@<$\langle$\textit{Check cmd args and Initialize n}$\rangle \thinspace +\!\!\equiv$@>
if (n < 0) {
    fprintf(stderr, "n >= 0\n");
    exit(1);
}
\end{oldcode}

Finally, we write the \textit{generating} portion of our code: \coderef{Generate Fibonacci sequence till n terms}.

\begin{code}[Generate Fibonacci sequence till n terms]
@<$\langle$\textit{Generate Fibonacci sequence till n terms}$\rangle\equiv$@>
for (int i = 0; i < n; ++i) {
    fprintf(stdout, "%d ", a + b);
    b = b + a;
    a = b - a;
}
\end{code}

%%%%%%%%%%%%%%%%%%%%%%%%%%%%%%%%%%%%%%%%%%%%%%%%%%%%%%%%%%%%%%%%%%%%%%%%%%%%%%%%%%%%%%%%%%%%%%%%%%%%%%%%%%%%%%
\subsection{Printing Prime Numbers}

In this example we are going to print prime numbers. Like in the previous example, we are going to begin by
first define the rules by which \tt prime.c \rm is to be built.

\begin{oldcode}
@<$\langle$\textit{Makefile}$\rangle \thinspace +\!\!\equiv$@>
CC := gcc -std=c99 -O0

runprimes: primes.exe
    ./primes.exe 30

primes.exe: primes.o
    $(CC) primes.o -o primes.exe

prime.o: prime.c
    $(CC) -c primes.c -o primes.o
\end{oldcode}

Now comes the interesting part: the C code which will print $n$ prime numbers. The core structure of our code
will be as follows:

\begin{code}[primes.c]
@<$\langle$\textit{primes.c}$\rangle\equiv$@>
#include <stdio.h>
#include <stdlib.h>

int isPrime(int);

int main(int argc, char **argv)
{
    @<\coderef{Initialize n}@>
    @<\coderef{Print prime numbers till n reaches 0}@>
    return 0;
}
\end{code}


Like in the previou example, $n$ will be supplied as a command line argument.

\begin{code}[Initialize n]
@<$\langle$\textit{Initialize n}$\rangle\equiv$@>
int n;
if (argc < 1) {
    fprintf(stderr, "n must be supplied as a command line argument\n");
    exit(1);
}
n = atoi(argv[1]);
\end{code}

We will have a loop in which n is decremented by 1 every time a prime number is encountered. Once it reaches
0, the loop ends.

\begin{code}[Print prime numbers till n reaches 0]
@<$\langle$\textit{Print prime numbers till n reaches 0}$\rangle\equiv$@>
int i = 0;
while (n > 0) {
    if (isPrime(i) == 0) {
        fprintf(stdout, "%d ", i);
        --n;
    }
    ++i;
}
fprintf(stdout, "Done\n");
\end{code}

In the loop, we are calling \tt isPrime() \rm to check whether the argument we have passed to it is a prime
number or not. This is done by finding the number of divisors of the argument from 2 till
$\frac{\text{argument}}{2}$. If the number of divisors is 0, then the argument is a prime number, else not.
The edge cases are for 0 and 1, which can be dealt with a simple \tt if \rm conditional statement.

\begin{code}[Check if i is prime or not]
@<$\langle$\textit{Check if i is prime or not}$\rangle\equiv$@>
int isPrime(int i)
{
    if (i == 0 || i == 1) {
        return 1;
    }
    for (int d = 2; d <= (i / 2); ++d) {
        if (i % d == 0) {
            return 1;
        }
    }
    return 0;
}
\end{code}

And we include it below \tt main() \rm in \coderef{primes.c}.

\begin{oldcode}
@<$\langle$\textit{primes.c}$\rangle \thinspace +\!\!\equiv$@>
@<\coderef{Check if i is prime or not}@>
\end{oldcode}

One more thing we should do is to write rules for cleaning the directory. 

\begin{oldcode}
@<$\langle$\textit{Makefile}$\rangle \thinspace +\!\!\equiv$@>

.PHONY: clean

clean: $(wildcard *.o) $(wildcard *.exe) $(wildcard *.out)
    rm -f $^
\end{oldcode}

%%%%%%%%%%%%%%%%%%%%%%%%%%%%%%%%%%%%%%%%%%%%%%%%%%%%%%%%%%%%%%%%%%%%%%%%%%%%%%%%%%%%%%%%%%%%%%%%%%%%%%%%%%%%%%
\subsection{Tangling and Weaving the Examples}

With our example codes now written, it is now time to tangle and weave them. Before tangling, we run \tt ldump
\rm to make a dump of the chunks in \texttt{example.web}. Then we can use \tt ltangle \rm to get out the
source code.

\begin{vcode}
$ ldump examples.web > examples.json
$ ltangle -i examples.json -R fibonacci.c -cc '//' > ../examples/fibonacci.c
$ ltangle -i examples.json -R primes.c -cc '//' > ../examples/primes.c
$ ltangle -i examples.json -R Makefile -cc '#' -ut 4 > ../examples/Makefile
\end{vcode}

To run this code, we can move to \tt examples/ \rm and run the \verb|make runfib| and \verb|make runprimes| to
print the first 30 terms of Fibonacci series and first 30 prime numbers, respectively.\\

To weave the documentation, simply run:
\begin{vcode}
$ lweave examples.web > documentation/examples.tex
\end{vcode}

Of course, you will need to have \tt litcode.sty \rm and a main container file from where you can
\verb|\input| the contents of \texttt{examples.tex}. The former can be arranged, by running \texttt{linit}.
For the latter, you will need to set up the container \texttt{tex} yourself.

\begin{vcode}
$ linit documentation/
\end{vcode}

